\documentclass[11pt,a4paper]{article}
\usepackage{xcolor}
\usepackage[utf8]{inputenc}
\usepackage{geometry}
\geometry{margin=1in}
\usepackage{hyperref}
\usepackage{enumitem}
\usepackage{listings}

\lstset{
  language=C,
  basicstyle=\ttfamily\small,
  keywordstyle=\bfseries,
  commentstyle=\itshape,
  showstringspaces=false,
  breaklines=true,
  captionpos=b,
  numbers=left,
  numberstyle=\tiny\color{gray},
}

\title{Digital Design Training Program \\ \large Day 6: The AVR Toolchain \& Matrix Keypad}
\author{Authored by }
\date{April 1, 2026}

\begin{document}

\maketitle

\section*{Theme: Transitioning away from automated build tools, writing custom hardware drivers in pure C, and implementing hardware interrupts for input logic.}

\subsection*{Module 6.1: Termux AVR-GCC \& Makefiles (45 Minutes)}

\begin{itemize}
    \item \textbf{Goal:} Understand and control the entire compilation process from C source code to a flashable \verb|.hex| file using the standard AVR-GCC toolchain and a \verb|Makefile|.
    \item \textbf{Action Steps:}
    \begin{enumerate}[label=\arabic*.]
        \item \textbf{Environment Setup:} Inside your Debian environment, install the essential tools.
        \begin{lstlisting}[language=bash]
# Update package lists and install the C compiler for AVR,
# binary utilities (for linking/objcopy), the standard C library for AVR,
# and the 'make' utility.
apt update
apt install gcc-avr binutils-avr avr-libc make
        \end{lstlisting}
        \item \textbf{The Compilation Pipeline Explained:}
        \begin{itemize}
            \item \textbf{Compiling:} \verb|avr-gcc| takes your \verb|.c| file and compiles it into an intermediate \verb|.o| (object) file containing machine code, but without memory addresses finalized.
            \item \textbf{Linking:} \verb|avr-gcc| then links all \verb|.o| files together into a single \verb|.elf| (Executable and Linkable Format) file. This file contains the final code, data, and metadata.
            \item \textbf{Extracting:} \verb|avr-objcopy| is a utility that extracts the raw program code from the \verb|.elf| file and formats it into the standard Intel \verb|.hex| format that programmers like ArduinoDroid understand.
        \end{itemize}
        \item \textbf{Writing a Makefile:} A \verb|Makefile| is a recipe for the \verb|make| utility. It defines how to build your project. Create a file named \verb|Makefile| with the following content:
\begin{lstlisting}[language=make]
# --- Project Definitions ---
# The microcontroller we are using
MCU = atmega328p
# The clock frequency of the microcontroller (16MHz for Arduino Uno)
F_CPU = 16000000UL
# The name of our project (will be the name of the .hex file)
TARGET = main

# --- Toolchain Definitions ---
CC = avr-gcc
OBJCOPY = avr-objcopy
CFLAGS = -Wall -Os -mmcu=$(MCU) -DF_CPU=$(F_CPU)

# --- Build Rules ---
# The default rule, executed when you just type ``make''
all: $(TARGET).hex

# Rule to create the .hex file from the .elf file
$(TARGET).hex: $(TARGET).elf
	$(OBJCOPY) -O ihex -R .eeprom $(TARGET).elf $(TARGET).hex

# Rule to link the .o file into a .elf file
$(TARGET).elf: $(TARGET).o
	$(CC) $(CFLAGS) -o $(TARGET).elf $(TARGET).o

# Rule to compile the .c file into a .o file
$(TARGET).o: $(TARGET).c
	$(CC) $(CFLAGS) -c $(TARGET).c -o $(TARGET).o

# Rule to clean up all generated files
clean:
	rm -f $(TARGET).o $(TARGET).elf $(TARGET).hex

# Rule to flash the code (assumes avrdude is installed and configured)
# This is a more advanced step, for now we will copy the .hex manually.
flash: $(TARGET).hex
	avrdude -c arduino -p $(MCU) -P /dev/ttyUSB0 -b 115200 -U flash:w:$(TARGET).hex
\end{lstlisting}
        \item \textbf{Execution:} Place your C code in a file named \verb|main.c| in the same directory as the \verb|Makefile|. Run \verb|make| in the terminal. It will automatically find the files and execute the commands to produce \verb|main.hex|.
    \end{enumerate}
    \item \textbf{Checkpoint:} You can successfully compile the bare-metal code from Day 5 using \verb|make| and flash the resulting \verb|.hex| file. You are now in full control of the build process.
\end{itemize}
\newpage
\subsection*{Module 6.2: Matrix Scanning in C (60 Minutes)}

\begin{itemize}
    \item \textbf{Goal:} Wire a 4x4 matrix keypad and write a custom bare-metal driver to read 16 buttons using only 8 microcontroller pins.
    \item \textbf{Conceptual Overview (Polling):} The keypad is a grid of wires. Pressing a button connects one row wire to one column wire. To find the pressed button, we can't just read the pins; we have to actively scan. The algorithm is:
    \begin{enumerate}
        \item Set all columns as outputs and all rows as inputs with pull-ups enabled.
        \item For each column:
        \begin{itemize}
            \item Drive only that column LOW.
            \item Read all the row pins.
            \item If a row pin is LOW, the button at the intersection of the active column and that row is the one being pressed.
        \end{itemize}
        \item Drive the column HIGH again and move to the next column.
    \end{enumerate}
    \item \textbf{Action Steps \& Code:}
\begin{lstlisting}[caption={Code for Module 6.2: Polling Keypad Scanner}]
#include <avr/io.h>
#include <util/delay.h>

// Let's use PORTC for Columns 0-3 and PORTD for Rows 0-3
#define COL_PORT PORTC
#define COL_DDR  DDRC
#define ROW_PORT PORTD
#define ROW_DDR  DDRD
#define ROW_PINS PIND

int main(void) {
    // Set LED on PB5 for output
    DDRB |= (1 << PB5);

    // Set rows as inputs with pull-ups
    ROW_DDR = 0x00;  // All pins of PORTD are inputs
    ROW_PORT = 0xFF; // Enable pull-ups on all pins of PORTD

    // Set columns as outputs, initially HIGH
    COL_DDR = 0xFF;  // All pins of PORTC are outputs
    COL_PORT = 0xFF; // All columns HIGH

    while (1) {
        int key_pressed = -1;

        // Scan through columns 0 to 3
        for (int c = 0; c < 4; c++) {
            // Pull the current column LOW
            COL_PORT &= ~(1 << c);

            // Check all rows
            for (int r = 0; r < 4; r++) {
                // If a row pin is now LOW, a button is pressed
                if (!(ROW_PINS & (1 << r))) {
                    key_pressed = (r * 4) + c;
                    // Simple debounce and wait for release
                    _delay_ms(50);
                    while(!(ROW_PINS & (1 << r)));
                    goto found; // Break out of both loops
                }
            }
            // Restore the column to HIGH before moving to the next
            COL_PORT |= (1 << c);
        }

    found:
        if (key_pressed != -1) {
            // Blink the LED 'key_pressed' number of times
            for(int i=0; i <= key_pressed; i++) {
                PORTB |= (1 << PB5); _delay_ms(100);
                PORTB &= ~(1 << PB5); _delay_ms(100);
            }
        }
    }
    return 0;
}
\end{lstlisting}
    \item \textbf{Checkpoint:} Pressing any of the 16 buttons correctly identifies its unique index (0-15) and blinks the onboard LED accordingly. The \verb|while(1)| loop is constantly busy scanning, consuming CPU cycles.
\end{itemize}
\newpage
\subsection*{Module 6.3: Pin Change Interrupts (PCINT) (60 Minutes)}

\begin{itemize}
    \item \textbf{Goal:} Eliminate wasteful polling by using hardware interrupts to detect keypresses, allowing the CPU to sleep or perform other tasks.
    \item \textbf{Conceptual Overview:} Instead of constantly asking ``Is a button pressed?'', we will configure the hardware to tell us when one is. We set the keypad to a ``rest state'' and enable Pin Change Interrupts on the row pins. When a button is pressed, a row pin changes state, triggering an ISR. The ISR's job is to then perform a single, quick scan to find out \emph{which} button was pressed.
    \item \textbf{Action Steps \& Code:}
\begin{lstlisting}[caption={Code for Module 6.3: Interrupt-Driven Keypad}]
#include <avr/io.h>
#include <avr/interrupt.h>
#include <util/delay.h>

// Using PORTD for rows, so we need PCINT on PORTD, which is PCINT2
#define ROW_PORT PORTD
#define ROW_DDR  DDRD
#define ROW_PINS PIND
#define COL_PORT PORTC
#define COL_DDR  DDRC

volatile int g_key = -1; // Global variable to store pressed key

// Function to put keypad pins in the ``listening'' state
void keypad_enable_interrupts() {
    // Rows are inputs with pull-ups
    ROW_DDR = 0x00;
    ROW_PORT = 0xFF;
    // Columns are outputs driven LOW
    COL_DDR = 0xFF;
    COL_PORT = 0x00;
    
    // Enable Pin Change Interrupt for PORTD (PCIE2)
    PCICR |= (1 << PCIE2);
    // Unmask interrupts for pins PD0-PD3 (PCINT16-19)
    PCMSK2 |= (1 << PCINT16) | (1 << PCINT17) | (1 << PCINT18) | (1 << PCINT19);
}

// The ISR for Pin Change Interrupt Request 2 (covers PORTD)
ISR(PCINT2_vect) {
    // Immediately disable the interrupt to prevent bouncing/re-entry
    PCICR &= ~(1 << PCIE2);

    // --- Perform a single scan to find the pressed key ---
    // Invert logic: Rows as outputs, Columns as inputs
    ROW_DDR = 0xFF;
    ROW_PORT = 0xFF;
    COL_DDR = 0x00;
    COL_PORT = 0xFF; // Enable pull-ups on columns
    _delay_us(10); // Wait for pull-ups to stabilize

    for (int r = 0; r < 4; r++) {
        ROW_PORT &= ~(1 << r); // Drive one row LOW
        for (int c = 0; c < 4; c++) {
            if (!(PINC & (1 << c))) { // Check if a column is now LOW
                g_key = (r * 4) + c;
                goto isr_found;
            }
        }
        ROW_PORT |= (1 << r); // Drive row HIGH again
    }

isr_found:
    // Store the key and re-enable interrupts for the next press
    keypad_enable_interrupts();
}

int main(void) {
    DDRB |= (1 << PB5); // LED for output
    keypad_enable_interrupts();
    sei(); // Enable global interrupts

    while (1) {
        // The main loop is now free! It only acts when a key has been found.
        if (g_key != -1) {
            cli(); // Disable interrupts to safely access g_key
            int key = g_key;
            g_key = -1;
            sei(); // Re-enable interrupts

            // Blink the LED to show the key was received
            for(int i=0; i <= key; i++) {
                PORTB |= (1 << PB5); _delay_ms(100);
                PORTB &= ~(1 << PB5); _delay_ms(100);
            }
        }
        // We could put the CPU to sleep here for max power savings!
    }
    return 0;
}
\end{lstlisting}
    \item \textbf{Checkpoint:} The keypad works as before, but the \verb|while(1)| loop is now empty. The CPU is only doing work when a button is physically pressed, which is a vastly more efficient design.
\end{itemize}

\end{document}
